% GENERATED FILE. Edit canonical lessons, then run npm run generate:books.
% canonical-source-hash: fnv1a64:154650005adbbdf9
% canonical-lessons: AR-C298-hariq, AR-C298-dukhan, AR-C298-ishtaala, AR-C298-ihtaraqa, AR-C298-rajul-itfa

\chapter{A fire}
\label{ch:ar-a-fire}
\hlchaptermodality{298}

\begin{chapteropening}
\textbf{By the end of this chapter:} \emph{I can say a blaze, smoke, to catch fire, to burn down and a firefighter.}

Complete the last lesson of chapter 298: I can say a blaze, smoke, to catch fire, to burn down and a firefighter.
\end{chapteropening}

\section[\texorpdfstring{\ar{حريق} (ḥarīq)}{حريق (ḥarīq)}]{\ar{حريق} (ḥarīq) — a blaze, a fire that spreads}

\label{lesson:AR-C298-hariq}



\begin{quote}
Before the new one: say the Arabic for to put on air, then the Arabic for on air.
\end{quote}

\subsection*{You'll want to know: \ar{حريق}}
\textbf{\ar{حريق}} — \emph{ḥarīq} — \textquotedblleft{}a blaze, a fire that spreads\textquotedblright{}.

\begin{quote}
\ar{رأى} \ar{الجيران} \ar{حريقا} \ar{في} \ar{المصنع}.
\end{quote}

\begin{quote}
\emph{raʾā al-jīrān ḥarīqan fī al-maṣnaʿ.}
\end{quote}

\begin{quote}
The neighbours saw a fire at the factory.
\end{quote}

\ar{نار} is fire itself; \ar{حريق} is a fire that destroys.

\begin{cousinweb}
\textbf{\ar{حريق}} (\emph{ḥarīq}) is built on the native Arabic root \emph{ḥ-r-q} (Lane), \textquotedblleft{}to burn\textquotedblright{}. \emph{Ḥarīq} is the fire that burns a thing up.
\end{cousinweb}

\begin{grammarlens}[title={What you've built}]
One more word for telling what happened.
\end{grammarlens}

\subsection*{Your turn}
\begin{itemize}
\raggedright
  \item \emph{Say it:} \emph{ḥarīq}
  \item \emph{Say it:} \emph{ḥarīq}, once more
  \item \emph{Say it:} \emph{raʾā al-jīrān ḥarīqan fī al-maṣnaʿ.} — The neighbours saw a fire at the factory.
  \item \emph{Recall:} say the Arabic for to put on air, then the Arabic for on air, then say \emph{ḥarīq} again
\end{itemize}

\begin{tcolorbox}[breakable,colback=teal!4,colframe=teal!35!black,title={Before you move on}]
What does \ar{حريق} mean? (\textquotedblleft{}A blaze, a fire that spreads\textquotedblright{}.) Say it once more.
\end{tcolorbox}
\section[\texorpdfstring{\ar{دخان} (dukhān)}{دخان (dukhān)}]{\ar{دخان} (dukhān) — smoke}

\label{lesson:AR-C298-dukhan}



\begin{quote}
Before the new one: say the Arabic for on air, then the Arabic for a blaze.
\end{quote}

\subsection*{You'll want to know: \ar{دخان}}
\textbf{\ar{دخان}} — \emph{dukhān} — \textquotedblleft{}smoke\textquotedblright{}.

\begin{quote}
\ar{خرج} \ar{دخان} \ar{أسود} \ar{من} \ar{النافذة}.
\end{quote}

\begin{quote}
\emph{kharaja dukhān aswad min an-nāfidha.}
\end{quote}

\begin{quote}
Black smoke came out of the window.
\end{quote}

In many dialects \ar{دخان} also means tobacco.

\begin{cousinweb}
\textbf{\ar{دخان}} (\emph{dukhān}) is built on the native Arabic root \emph{d-kh-n} (Lane), where \emph{dakhana} is said of a fire that sends up smoke.
\end{cousinweb}

\begin{grammarlens}[title={What you've built}]
One more word for telling what happened.
\end{grammarlens}

\subsection*{Your turn}
\begin{itemize}
\raggedright
  \item \emph{Say it:} \emph{dukhān}
  \item \emph{Say it:} \emph{dukhān}, once more
  \item \emph{Say it:} \emph{kharaja dukhān aswad min an-nāfidha.} — Black smoke came out of the window.
  \item \emph{Recall:} say the Arabic for on air, then the Arabic for a blaze, then say \emph{dukhān} again
\end{itemize}

\begin{tcolorbox}[breakable,colback=teal!4,colframe=teal!35!black,title={Before you move on}]
What does \ar{دخان} mean? (\textquotedblleft{}Smoke\textquotedblright{}.) Say it once more.
\end{tcolorbox}
\section[\texorpdfstring{\ar{اشتعل} (ishtaʿala)}{اشتعل (ishtaʿala)}]{\ar{اشتعل} (ishtaʿala) — to catch fire, to blaze}

\label{lesson:AR-C298-ishtaala}



\begin{quote}
Before the new one: say the Arabic for a blaze, then the Arabic for smoke.
\end{quote}

\subsection*{You'll want to know: \ar{اشتعل}}
\textbf{\ar{اشتعل}} — \emph{ishtaʿala} — \textquotedblleft{}to catch fire, to blaze\textquotedblright{}.

\begin{quote}
\ar{اشتعلت} \ar{النار} \ar{في} \ar{الخشب} \ar{بسرعة}.
\end{quote}

\begin{quote}
\emph{ishtaʿalat an-nār fī al-khashab bi-surʿa.}
\end{quote}

\begin{quote}
The fire caught in the wood quickly.
\end{quote}

Present: \ar{يشتعل}, \emph{yashtaʿil}.

\begin{cousinweb}
\textbf{\ar{اشتعل}} (\emph{ishtaʿala}) is built on the native Arabic root \emph{sh-ʿ-l} (Lane), \textquotedblleft{}to kindle a fire\textquotedblright{}. It is Form VIII, \textquotedblleft{}to take fire\textquotedblright{}; \emph{shuʿla} is a flame.
\end{cousinweb}

\begin{grammarlens}[title={What you've built}]
One more word for telling what happened.
\end{grammarlens}

\subsection*{Your turn}
\begin{itemize}
\raggedright
  \item \emph{Say it:} \emph{ishtaʿala}
  \item \emph{Say it:} \emph{ishtaʿala}, once more
  \item \emph{Say it:} \emph{ishtaʿalat an-nār fī al-khashab bi-surʿa.} — The fire caught in the wood quickly.
  \item \emph{Recall:} say the Arabic for a blaze, then the Arabic for smoke, then say \emph{ishtaʿala} again
\end{itemize}

\begin{tcolorbox}[breakable,colback=teal!4,colframe=teal!35!black,title={Before you move on}]
What does \ar{اشتعل} mean? (\textquotedblleft{}To catch fire, to blaze\textquotedblright{}.) Say it once more.
\end{tcolorbox}
\section[\texorpdfstring{\ar{احترق} (iḥtaraqa)}{احترق (iḥtaraqa)}]{\ar{احترق} (iḥtaraqa) — to burn down, to get burnt}

\label{lesson:AR-C298-ihtaraqa}



\begin{quote}
Before the new one: say the Arabic for smoke, then the Arabic for to catch fire.
\end{quote}

\subsection*{You'll want to know: \ar{احترق}}
\textbf{\ar{احترق}} — \emph{iḥtaraqa} — \textquotedblleft{}to burn down, to get burnt\textquotedblright{}.

\begin{quote}
\ar{احترق} \ar{الخبز} \ar{في} \ar{الفرن}.
\end{quote}

\begin{quote}
\emph{iḥtaraqa al-khubz fī al-furn.}
\end{quote}

\begin{quote}
The bread got burnt in the oven.
\end{quote}

Present: \ar{يحترق}, \emph{yaḥtariq}.

\begin{cousinweb}
\textbf{\ar{احترق}} (\emph{iḥtaraqa}) is built on the native Arabic root \emph{ḥ-r-q} (Lane), \textquotedblleft{}to burn\textquotedblright{}. It is Form VIII, \textquotedblleft{}to be burnt, to burn up\textquotedblright{}.
\end{cousinweb}

\begin{grammarlens}[title={What you've built}]
One more word for telling what happened.
\end{grammarlens}

\subsection*{Your turn}
\begin{itemize}
\raggedright
  \item \emph{Say it:} \emph{iḥtaraqa}
  \item \emph{Say it:} \emph{iḥtaraqa}, once more
  \item \emph{Say it:} \emph{iḥtaraqa al-khubz fī al-furn.} — The bread got burnt in the oven.
  \item \emph{Recall:} say the Arabic for smoke, then the Arabic for to catch fire, then say \emph{iḥtaraqa} again
\end{itemize}

\begin{tcolorbox}[breakable,colback=teal!4,colframe=teal!35!black,title={Before you move on}]
What does \ar{احترق} mean? (\textquotedblleft{}To burn down, to get burnt\textquotedblright{}.) Say it once more.
\end{tcolorbox}
\section[\texorpdfstring{\ar{رجل} \ar{إطفاء} (rajul iṭfāʾ)}{رجل إطفاء (rajul iṭfāʾ)}]{\ar{رجل} \ar{إطفاء} (rajul iṭfāʾ) — a firefighter}

\label{lesson:AR-C298-rajul-itfa}



\begin{quote}
Before the new one: say the Arabic for to catch fire, then the Arabic for to burn down.
\end{quote}

\subsection*{You'll want to know: \ar{رجل} \ar{إطفاء}}
\textbf{\ar{رجل} \ar{إطفاء}} — \emph{rajul iṭfāʾ} — \textquotedblleft{}a firefighter\textquotedblright{}.

\begin{quote}
\ar{صعد} \ar{رجل} \ar{الإطفاء} \ar{إلى} \ar{السطح}.
\end{quote}

\begin{quote}
\emph{ṣaʿida rajul al-iṭfāʾ ilā as-saṭḥ.}
\end{quote}

\begin{quote}
The firefighter climbed up to the roof.
\end{quote}

Plural \ar{رجال} \ar{الإطفاء}. The fire brigade is \ar{الإطفاء}.

\begin{cousinweb}
\textbf{\ar{رجل} \ar{إطفاء}} (\emph{rajul iṭfāʾ}) is \textquotedblleft{}a man of extinguishing\textquotedblright{}: \emph{iṭfāʾ} is the verbal noun of the taught \emph{aṭfaʾa}, built on the native Arabic root \emph{ṭ-f-ʾ}, \textquotedblleft{}(of fire) to go out\textquotedblright{} (Lane).
\end{cousinweb}

\begin{grammarlens}[title={What you've built}]
One more word for telling what happened.
\end{grammarlens}

\subsection*{Your turn}
\begin{itemize}
\raggedright
  \item \emph{Say it:} \emph{rajul iṭfāʾ}
  \item \emph{Say it:} \emph{rajul iṭfāʾ}, once more
  \item \emph{Say it:} \emph{ṣaʿida rajul al-iṭfāʾ ilā as-saṭḥ.} — The firefighter climbed up to the roof.
  \item \emph{Recall:} say the Arabic for to catch fire, then the Arabic for to burn down, then say \emph{rajul iṭfāʾ} again
\end{itemize}

\begin{tcolorbox}[breakable,colback=teal!4,colframe=teal!35!black,title={Before you move on}]
What does \ar{رجل} \ar{إطفاء} mean? (\textquotedblleft{}A firefighter\textquotedblright{}.) Say it once more.
\end{tcolorbox}
